\begin{helper}
Let $n\ge1$ be an integer, let $p>1$ be real, let $A\colon\mathbb{Z}\to\mathbb{C}$
satisfy
\[
  \sum_{m=0}^\infty(n+m)\,|A_{n+m}|^2<\infty ,
\]
write $A_k=A(k)$, $\lambda_k=\lambda_k(n,p)$ for the exponent of
\noderef{LambdaCoef}, and let $T_k=T_k(n,p,A,\cdot)$ be the $k$-th summand of
the series (19) as in \noderef{Thm2Term}.  Then the function
\[
  H(\xi)=\begin{cases}A_n,&\xi=0,\\[2pt]
    \displaystyle\sum_{m=0}^\infty T_{n+m}(\xi),&\xi\ne0,\end{cases}
\]
is continuous on the whole unit disc $\mathbb{B}=\{\xi:|\xi|<1\}$, the origin
included.

This is conjunct (a) of Theorem~2 for the constructed solution, the statement
$H\in C(\mathbb{B})$ of \texttt{paper.tex:264}.  Off the origin it is local
uniform convergence of a series of continuous summands.  At the origin it is
the real content, and it is what forces the completing value: the $m=0$
summand of (19) is the constant $A_n$, every later summand carries a positive
power of $|\xi|$, and so the sum tends to $A_n$ as $\xi\to0$.  No other value
at the origin would make the function continuous there.

The sum is written out as a case distinction rather than named, so that the
statement depends on nothing but $n$, $p$ and $A$.
\end{helper}
\begin{proof}
By \noderef{Thm2TermGeometricSummability}, applied with the hypothesis on $A$,
fix once and for all a rate $b$ with $0<b\le\tfrac12$ and a bound $CA\ge0$ such
that
\begin{equation}\tag{1}
  |A_{n+m}|\le CA\quad\text{for all }m,
  \qquad
  |\xi|^{\lambda_{n+m}}\le\bigl(|\xi|^{b}\bigr)^{m}
    \quad\text{for }0<|\xi|<1\text{ and all }m,
\end{equation}
and such that, for $0<\rho<1$,
\begin{equation}\tag{2}
  \|T_{n+m}(\xi)\|\le2\,|A_{n+m}|\,|\xi|^{\lambda_{n+m}}\ (\xi\ne0),
  \qquad
  \sum_{m=0}^\infty2\,|A_{n+m}|\,\rho^{\lambda_{n+m}}<\infty .
\end{equation}
The same node gives $T_n(\xi)=A_n$ for $\xi\ne0$, the $m=0$ summand being
constant.  Finally $\lambda_{n+m}\ge0$ for every $m$, by
\noderef{LambdaEpsEstimates}: $\lambda_n=0$, and for $m\ge1$ one has
$\lambda_{n+m}\ge(p-1)/(np)>0$.  Each summand $T_{n+m}$ is continuous away
from the origin, by part (a) of \noderef{Thm2TermSolvesEq18} applied to the
index $n+m\ge n$.

Two preliminary remarks.  First, for $\xi\ne0$ with $|\xi|<1$ the series
$\sum_m\|T_{n+m}(\xi)\|$ converges: by (2) its terms are at most
$2|A_{n+m}|\,|\xi|^{\lambda_{n+m}}$, which is summable by the second half of
(2) at $\rho=|\xi|$.  So the defining sum exists at every such point, and
$H(\xi)=\sum_mT_{n+m}(\xi)$ there.  Second, $H(0)=A_n$ by definition.

\emph{Step 1: continuity at a point $x\ne0$ of $\mathbb{B}$.}  Put
$\delta=|x|/2$ and $\rho=(|x|+1)/2$, so that
\[
  0<\delta<|x|<\rho<1 ,
\]
and let
\[
  U=\bigl\{z:\delta<|z|<\rho\bigr\},
\]
an open annulus, being the intersection of the two open sets $\{|z|>\delta\}$
and $\{|z|<\rho\}$; it contains $x$.  Every $z\in U$ is nonzero, since
$|z|>\delta>0$.

On $U$ each summand $T_{n+m}$ is continuous, by part (a) of
\noderef{Thm2TermSolvesEq18} as recalled above, and is bounded by a constant
independent of $z$:
\[
  \|T_{n+m}(z)\|\le2\,|A_{n+m}|\,|z|^{\lambda_{n+m}}
  \le2\,|A_{n+m}|\,\rho^{\lambda_{n+m}},
\]
the first inequality by (2) and the second because $|z|<\rho$ and
$\lambda_{n+m}\ge0$, so that $t\mapsto t^{\lambda_{n+m}}$ is nondecreasing on
$[0,\infty)$.  The bounding series $\sum_m2|A_{n+m}|\rho^{\lambda_{n+m}}$
converges by (2), since $0<\rho<1$.  Hence
$\sum_mT_{n+m}$ converges uniformly on $U$ and its sum
\[
  G(z)=\sum_{m=0}^\infty T_{n+m}(z)
\]
is continuous on $U$, by the Weierstrass criterion for a uniformly convergent
series of continuous functions.

On $U$ we have $H=G$, because every point of $U$ is nonzero.  Therefore $H$ is
continuous on $U$, and $U$ is an open neighbourhood of $x$, so $H$ is
continuous at $x$; in particular continuous at $x$ within $\mathbb{B}$.

\emph{Step 2: the tail estimate near the origin.}  Let $\xi\ne0$ with
$|\xi|<1$ and suppose moreover $|\xi|^{b}\le\tfrac12$.  Write
$q=|\xi|^{b}$, so $0<q\le\tfrac12$.  Splitting off the $m=0$ term of the
(absolutely convergent) defining series and using $T_n(\xi)=A_n$,
\begin{equation}\tag{3}
  H(\xi)-A_n=\sum_{m=1}^\infty T_{n+m}(\xi).
\end{equation}
For $m\ge1$, combining (2), the second half of (1) and $|A_{n+m}|\le CA$,
\[
  \|T_{n+m}(\xi)\|
  \le2\,|A_{n+m}|\,|\xi|^{\lambda_{n+m}}
  \le2\,CA\,q^{m}.
\]
Since $m\ge1$ and $0<2q\le1$ we have $(2q)^{m}\le2q$, i.e.
$q^{m}\le2q\cdot2^{-m}$, whence
\begin{equation}\tag{4}
  \|T_{n+m}(\xi)\|\le4\,CA\,q\cdot\Bigl(\frac12\Bigr)^{m}
  \qquad(m\ge1).
\end{equation}
The right-hand side of (4) is summable in $m$, being a constant times a
geometric series of ratio $\tfrac12$, so by (3), the triangle inequality for
absolutely convergent series and (4),
\begin{equation}\tag{5}
  \|H(\xi)-A_n\|
  \le\sum_{m=1}^\infty4\,CA\,q\Bigl(\frac12\Bigr)^{m}
  \le4\,CA\,q\cdot\frac1{1-\frac12}
  =8\,CA\,|\xi|^{b}.
\end{equation}
(The estimate (5) is stated for $\xi\ne0$; at $\xi=0$ its left-hand side is
$0$ by $H(0)=A_n$, so it holds there as well.)

\emph{Step 3: continuity at the origin.}  Since $b>0$, the function
$t\mapsto t^{b}$ is continuous at $t=0$ with value $0^{b}=0$, and $\xi\mapsto|\xi|$
is continuous with $|0|=0$; composing, $|\xi|^{b}\to0$ as $\xi\to0$, hence also
as $\xi\to0$ within $\mathbb{B}$.  Consequently $8\,CA\,|\xi|^{b}\to0$ along
the same filter, and the condition $|\xi|^{b}\le\tfrac12$ holds eventually
along it.

So, eventually as $\xi\to0$ within $\mathbb{B}$, the hypotheses of Step 2 are
met and $\|H(\xi)-A_n\|\le8\,CA\,|\xi|^{b}$ by (5); the dominating quantity
tends to $0$; therefore $H(\xi)\to A_n=H(0)$ as $\xi\to0$ within
$\mathbb{B}$.  That is continuity of $H$ at $0$ within $\mathbb{B}$.

\emph{Step 4: conclusion.}  Let $x\in\mathbb{B}$.  If $x=0$, $H$ is continuous
at $x$ within $\mathbb{B}$ by Step 3.  If $x\ne0$, it is continuous at $x$ by
Step 1, hence continuous at $x$ within $\mathbb{B}$.  As $x$ was arbitrary, $H$
is continuous on $\mathbb{B}$.
\end{proof}
